\documentclass[10pt,a4paper]{amsart}
\usepackage[margin=19mm]{geometry}
\usepackage{amsmath,amssymb,mathtools}
\usepackage[scr=rsfs,cal=euler]{mathalfa}
\usepackage{xfrac}
\usepackage[unicode,hidelinks,bookmarks=false]{hyperref}
\newcommand{\ie}{i.e.\ }
\newcommand{\cf}{cf.\ }
\newcommand{\resp}{respectively\ }
\newcommand{\Subsection}{\subsection}
\providecommand{\nobreakdash}{\nobreak\hbox{-}}
\setlength{\emergencystretch}{2em}
\multlinegap0pt
\usepackage{bm}
\usepackage{bbm}
\usepackage{dsfont}
\usepackage{xspace}
\usepackage{cancel}
\usepackage{mathtools}
\usepackage{amsthm}
\makeatletter
\@ifundefined{theorem}{\newtheorem{theorem}{Theorem}}{}
\makeatother
\title[The Haar State values of monomials and a method to pick orth]{The Haar State values of monomials and a method to pick orthonormal bases on $O(U\_q(3))$}
\author{Ting Lu}
\date{Source: arXiv:2508.15346v1 (CC BY 4.0)}
\begin{document}
\maketitle

\noindent\emph{Source and licence.} The Haar State values of monomials and a method to pick orthonormal bases on $O(U_q(3))$. Version arXiv:2508.15346v1 (CC BY 4.0), distributed under the Creative Commons Attribution 4.0 International licence, as recorded in the arXiv licence metadata for this exact version and shown on the abstract page https://arxiv.org/abs/2508.15346v1. Full rights notice: licenses/arxiv-2508.15346-CC-BY-4.0.txt.\par\smallskip

\noindent\textit{Editorial scope.} Counted unit: the Theorem whose source label is \texttt{eq:4.73} in the article (source0.tex lines 1685--1699, source label \texttt{eq:4.73}), delivered together with the article's own complete original proof of it (source0.tex lines 1704--1803, one \texttt{proof} environment of 100 lines). No mathematics is written, repaired or re-derived: statement and proof are the article's own text line for line. The proof is detached from the statement in the source: the article states the result at lines 1685--1699 and prints its proof at lines 1704--1803, and the proof environment's own header names the statement, so the association is the article's own. Required context retained uncounted: eq:4.57 (lines 1443--1458); eq:4.63 (lines 1571--1576); eq:4.69 (lines 1626--1638); eq:4.72 (lines 1675--1683). Every definition, display and label of those lines is retained unchanged. Omitted from this package and not redistributed: 1--1442, 1459--1570, 1577--1625, 1639--1674, 1684--1684, 1700--1703, 1804--3031. Nothing inside a retained excerpt is omitted or altered: every retained body is delivered exactly as the article's lines are written, so no omission receipt of any kind accompanies this package. Citations: the retained excerpts carry no citation command at all, so no bibliography entry of the article is transcribed and no reference is redistributed. Reference adaptation: none was needed and none was made; every reference inside the delivered unit resolves inside it, so no unresolved reference or citation is produced. Printed slips kept literal: no printed word, symbol, hypothesis or number of the counted statement or of its proof was altered. Layout and class: the article's own class is \texttt{amsart} with packages including graphicx, adjustbox, comment, amsmath, amsthm, amssymb, mathrsfs, diagbox, mathtools, tikz, pgfplots, rotating, lscape, multicol; the wrapper is the lane's pinned \texttt{amsart} editorial wrapper at 10pt on A4, which restates the theorem-like environments the retained text needs and adds the article's own macro definitions that the retained text needs (none), with extra wrapper packages (bm, bbm, dsfont, xspace, cancel, mathtools). The delivered numbering is the unit's own. Assets: no retained span contains a figure environment, a graphics-inclusion command, a PDF-page inclusion, a table or a TikZ picture; no image file is copied into this package, the package declares no asset dependency and no image file is redistributed with it. Licence: version arXiv:2508.15346v1 of this article is distributed under the Creative Commons Attribution 4.0 International licence, as recorded in the arXiv licence metadata for this exact version and shown on the abstract page https://arxiv.org/abs/2508.15346v1. A full rights notice is delivered at licenses/arxiv-2508.15346-CC-BY-4.0.txt. Characters: every retained excerpt is pure ASCII, so the delivered engine, XeLaTeX with its default Latin Modern text font, reports no missing character.
\begin{equation}\label{eq:4.57}
    \begin{split}
        &h(a^sb^{m-r-s}c^{r}d^{m-s-l}e^{l+r+s-m}f^{m-r}g^{l}h^{m-l}\cdot\text{det}_q^{-m})\\
        % =&(-1)^{l+r+s-m}q^{(m-2r+l+s+1)(l+r+s-m)}\frac{\prod_{i=m-s-l+1}^{r}(1-q^{2i})}{\prod_{j=m-r+1}^{l+s}(1-q^{2j})}\\
        % &\frac{(-1)^{m-s}q^{(2m+1)(m-s)-2m-2l^2-2(m-s-l)^2-(m-s-l)l}}{{m\choose m-s-l}_{q^2}{m\choose l}_{q^2}}\\
        % &h\left(c^me^mg^m\cdot\text{det}_q^{-m}\right)\\
        %(-2r+l+2s+1)(l+r+s-m)+(2m+1)(m-s)-2(m-s-l)^2-(m-s-l)l\\
        %(-2r+2l+2s+1)(l+r+s-m)+(2m+1)(m-s)-2(m-s-l)^2-rl\\
        %(-2r+2l+2s+1)r+(-2r+1+2m))(l+s-m)+(2m+1)(m-s)-rl\\
        %(-2r+2l+2s+1)r+l(-2r+1+2m)-2r(s-m)-rl\\
        %-2r^2+2lr+2sr+r-2rl+l+2ml-2r(s-m)-rl\\
        %r+l+2ml+2rm-rl-2r^2\\
        =&\frac{(-1)^{r+l}q^{(2m+1)(l+r)-2l^2-2r^2-rl-2m+(m-s)(l+r+s-m)}}{{m\choose r}_{q^2}{m\choose l}_{q^2}}\\
        &\cdot h\left(c^me^mg^m\cdot\text{det}_q^{-m}\right).
    \end{split}
\end{equation}

\begin{equation}\label{eq:4.63}
    \begin{split}
        &h\left(a^sb^{m-s-r}c^rd^{m-s-l}f^{m-r-t}g^lh^{m-l-t}k^{t}\cdot\text{det}_q^{-m}\right)\\
        =&\frac{(-1)^{r+l}q^{4st+(2m+1)(l+r)-2m-2l^2-2r^2-rl}}{{m\choose r+t}_{q^2}{m\choose l+t}_{q^2}}h\left(c^me^mg^m\cdot\text{det}_q^{-m}\right).
    \end{split}
\end{equation}

\begin{equation}\label{eq:4.69}
    \begin{split}
        &h\left(a^{s}b^{m-s}d^{m-s-l}e^{s+l+t-m}f^{m-t}g^{l}h^{m-l-t}k^{t}\cdot\text{det}_q^{-m}\right)\\
        % =&(-1)^{m-s-t}q^{(s+l+t-m)(-2t-1+2s+l)}\frac{\prod_{i=m-s-l+1}^{t}(1-q^{2j})}{\prod_{i=m-t}^{s+l}(1-q^{2i})}\\
        % &\frac{q^{4s(m-s-l)+(2m+1)l-2m-2l^2}}{{m\choose m-s-l}_{q^2}{m\choose m-s}_{q^2}}h\left(c^me^mg^m\cdot\text{det}_q^{-m}\right)\\
        % =&(-1)^{m-s-t}
        % \frac{q^{4st-2(s+l+t-m)^2-(2m+1-3l)(s+l+t-m)+(2m+1)l-2m-2l^2}}{{m\choose t}_{q^2}{m\choose s}_{q^2}}\\
        % &\cdot h\left(c^me^mg^m\cdot\text{det}_q^{-m}\right)\\
        =&(-1)^{m-s-t}
        \frac{q^{4st-(2s+2t+1-l)(s+l+t-m)+(2m+1)l-2m-2l^2}}{{m\choose t}_{q^2}{m\choose s}_{q^2}}\\
        &\cdot h\left(c^me^mg^m\cdot\text{det}_q^{-m}\right).
    \end{split}
\end{equation}

\begin{equation}\label{eq:4.72}
    \begin{split}
        &h\left(a^sb^{m-s-r}c^rd^{m-s-l}e^{n}f^{m-r-t}g^lh^{m-l-t}k^t\cdot\text{det}_q^{-m}\right)\\
        =&-q^{3m-s-l-4r-3t+3}\frac{1-q^{2r}}{1-q^{2(m-r-t+1)}}\\
        &\cdot h\left(a^sb^{m-s-r+1}c^{r-1}d^{m-s-l}e^{n-1}f^{m-r-t+1}g^lh^{m-l-t}k^t\cdot\text{det}_q^{-m}\right)\\
        &-q^{2m-l-4t-r+1}\frac{1-q^{2t}}{1-q^{2(m-r-t+1)}}\\
        &\cdot h\left(a^sb^{m-s-r}c^rd^{m-s-l}e^{n-1}f^{m-r-t+1}g^lh^{m-l-t+1}k^{t-1}\cdot\text{det}_q^{-m}\right).
    \end{split}
\end{equation}

\begin{theorem}[Main Theorem]
     Let $m,s,r,l,t$ be a set of integers satisfying condition (4.3.a-c) and set $n=s+r+l+t-m$. We have
     \begin{equation}\tag{4.23}
         h\left(c^me^mg^m\cdot\text{det}_q^{-m}\right)=\frac{(-q)^{3m}(q^2-1)^2(q^4-1)}{(q^{2m+2}-1)^2(q^{2m+4}-1)},
     \end{equation}
     and
    \begin{equation}\label{eq:4.73}
        \begin{split}
            &h(a^sb^{m-s-r}c^rd^{m-s-l}e^{n}f^{m-r-t}g^lh^{m-l-t}k^t\cdot\text{det}_q^{-m})\\
            =&\left(\sum_{k=0}^{n}(-1)^{k}q^{(n-k)(n-3k-1)+2k(s+t)}\frac{{r\choose k}_{q^2}{l\choose k}_{q^2}{s\choose n-k}_{q^2}{t\choose n-k}_{q^2}}{{n\choose k}_{q^2}}\right)\\
            &\frac{(-1)^{r+l+n}q^{(2m+1)(l+r)+(n-2)m-2l^2-2r^2-rl+4st-3ns-3nt}}{{m\choose n,m-l-s,m-r-t}_{q^2}{m\choose n,m-r-s,m-l-t}_{q^2}}\\
            &\cdot h\left(c^me^mg^m\cdot\text{det}_q^{-m}\right).
        \end{split}
    \end{equation}
\end{theorem}

\begin{proof}[Proof of (\ref{eq:4.73})]
    It is easy to check that our initial condition (\ref{eq:4.63}) and boundary condition (\ref{eq:4.57}) and (\ref{eq:4.69}) satisfies (\ref{eq:4.73}). 
    % The initial condition corresponds to $n=0$. It is to check that (\ref{eq:4.73}) reduce to (\ref{eq:4.63}) when $n=0$. The boundary condition (\ref{eq:4.57}) corresponds to $t=0$. Thus, terms in the parenthesis of (\ref{eq:4.73}) are non-zero only if $k=n$ and (\ref{eq:4.73}) reduces to
    % \begin{equation}
    %     \begin{split}
    %         &q^{2ns}{r\choose n}_{q^2}{l\choose n}_{q^2}\frac{(-1)^{r+l}q^{(2m+1)(l+r)+(n-2)m-2l^2-2r^2-rl-3ns}}{{m\choose n,r-n,m-r}_{q^2}{m\choose n,l-n,m-l}_{q^2}}\\
    %         &\cdot h\left(c^me^mg^m\cdot\text{det}_q^{-m}\right)\\
    %         =&\frac{(-1)^{r+l}q^{(2m+1)(l+r)-2m-2l^2-2r^2-rl+n(m-s)}}{{m\choose r}_{q^2}{m\choose l}_{q^2}}h\left(c^me^mg^m\cdot\text{det}_q^{-m}\right)
    %     \end{split}
    % \end{equation}
    % which is consistent with (\ref{eq:4.57}). The boundary condition (\ref{eq:4.69}) corresponds to $r=0$. Thus, terms in the parenthesis of (\ref{eq:4.73}) is non-zero only if $k=0$. Thus, (\ref{eq:4.73}) reduces to
    % \begin{equation}
    %     \begin{split}
    %         &q^{n(n-1)}{s\choose n}_{q^2}{t\choose n}_{q^2}\frac{(-1)^{l+n}q^{(2m+1)l+(n-2)m-2l^2+4st-3ns-3nt}}{{m\choose n,t-n,m-t}_{q^2}{m\choose n,m-s,s-n}_{q^2}}\\
    %         &\cdot h\left(c^me^mg^m\cdot\text{det}_q^{-m}\right)\\
    %         =&q^{n(n-1)}\frac{(-1)^{l+n}q^{(2m+1)l+(n-2)m-2l^2+4st-3ns-3nt}}{{m\choose t}_{q^2}{m\choose s}_{q^2}}\\
    %         &\cdot h\left(c^me^mg^m\cdot\text{det}_q^{-m}\right)\\
    %         =&\frac{(-1)^{m-s-t}q^{(2m+1)l-2m-2l^2+4st-n(2s+2t+1-l)}}{{m\choose t}_{q^2}{m\choose s}_{q^2}}h\left(c^me^mg^m\cdot\text{det}_q^{-m}\right)
    %     \end{split}
    % \end{equation}
    % which is consistent with (\ref{eq:4.69}).

    In the following computation, we will omit the part $h\left(c^me^mg^m\cdot\text{det}_q^{-m}\right)$ for simplicity. Recall the following expression from the theory of basic hypergeometric functions
    \begin{equation}\label{eq:4.71+}
        (a;q)_n=\prod_{i=0}^{n-1}(1-aq^i) \text{ and }(a;q)_0=1.
    \end{equation}
    With this notation we can rewrite (\ref{eq:4.73}) as
    \begin{equation}\label{eq:4.77}
        \begin{split}
            &h\left(a^sb^{m-s-r}c^rd^{m-s-l}e^{n}f^{m-r-t}g^lh^{m-l-t}k^t\cdot\text{det}_q^{-m}\right)\\
            =&\left(\sum_{k=0}^{n}(-1)^{k}q^{(n-k)(n-3k-1)+2k(s+t)}\right.\\
            &\left.\cdot{n\choose k}_{q^{2}}(q^{2r};q^{-2})_k(q^{2l};q^{-2})_k(q^{2s};q^{-2})_{n-k}(q^{2t};q^{-2})_{n-k}\right)\\
            &\frac{(q^2;q^2)_{r+t-n}(q^2;q^2)_{l+t-n}(q^2;q^2)_{r+s-n}(q^2;q^2)_{l+s-n}}{(q^2;q^2)_m^2}\\
            &(-1)^{r+l+n}q^{(2m+1)(l+r)+(n-2)m-2l^2-2r^2-rl+4st-3ns-3nt}.
        \end{split}
    \end{equation}
    Now, assume that (\ref{eq:4.77}) is true for all $n\le p-1$. We show that (\ref{eq:4.77}) also hold for $n=p$ using the recursive relation (\ref{eq:4.72}). Substituting the expression (\ref{eq:4.77}) into the $h(m;s,r-1,l,t)$ and $h(m;s,r,l,t-1)$ in (\ref{eq:4.72}) and after simplification, we get
    % \begin{equation}
    %     \begin{split}
    %         &h\left(a^sb^{m-s-r}c^rd^{m-s-l}e^{p}f^{m-r-t}g^lh^{m-l-t}k^t\cdot\text{det}_q^{-m}\right)\\
    %         =&-q^{3m-s-l-4r-3t+3}\frac{1-q^{2r}}{1-q^{2(m-r-t+1)}}\\
    %         &\left(\sum_{k=0}^{p-1}(-1)^{k}q^{(p-1-k)(p-1-3k-1)+2k(s+t)}\right.\\
    %         &\left.\cdot{p-1\choose k}_{q^{2}}(q^{2r-2};q^{-2})_k(q^{2l};q^{-2})_k(q^{2s};q^{-2})_{p-1-k}(q^{2t};q^{-2})_{p-1-k}\right)\\
    %         &\frac{(q^2;q^2)_{r+t-p}(q^2;q^2)_{l+t-p+1}(q^2;q^2)_{r+s-p}(q^2;q^2)_{l+s-p+1}}{(q^2;q^2)_m^2}\\
    %         &(-1)^{r+l+p}q^{(2m+1)(l+r-1)+(p-3)m-2l^2-2(r-1)^2-(r-1)l+4st-3(p-1)s-3(p-1)t}\\
    %         &-q^{2m-l-4t-r+1}\frac{1-q^{2t}}{1-q^{2(m-r-t+1)}}\\
    %         &\left(\sum_{k=0}^{p-1}(-1)^{k}q^{(p-1-k)(p-1-3k-1)+2k(s+t-1)}\right.\\
    %         &\left.\cdot{p-1\choose k}_{q^{2}}(q^{2r};q^{-2})_k(q^{2l};q^{-2})_k(q^{2s};q^{-2})_{p-1-k}(q^{2t-2};q^{-2})_{p-1-k}\right)\\
    %         &\frac{(q^2;q^2)_{r+t-p}(q^2;q^2)_{l+t-p}(q^2;q^2)_{r+s-p+1}(q^2;q^2)_{l+s-p+1}}{(q^2;q^2)_m^2}\\
    %         &(-1)^{r+l+p-1}q^{(2m+1)(l+r)+(p-3)m-2l^2-2r^2-rl+4st-3(p-1)s-3(p-1)(t-1)}.
    %     \end{split}
    % \end{equation}
    \begin{equation}\label{eq:4.78}
        \begin{split}
            &h\left(a^sb^{m-s-r}c^rd^{m-s-l}e^{p}f^{m-r-t}g^lh^{m-l-t}k^t\cdot\text{det}_q^{-m}\right)\\
            =&\frac{(q^2;q^2)_{r+t-p}(q^2;q^2)_{l+t-p}(q^2;q^2)_{r+s-p}(q^2;q^2)_{l+s-p}}{(q^2;q^2)_m^2}\\
            &(-1)^{r+l+p}q^{(2m+1)(l+r)+(p-2)m-2l^2-2r^2-rl+4st-3ps-3pt}\\
            &\left(\sum_{k=0}^{p-1}(-1)^{k}q^{(p-1-k)(p-3k-2)+2k(s+t)}\right.\\
            &\cdot{p-1\choose k}_{q^{2}}(q^{2l};q^{-2})_k(q^{2r};q^{-2})_k(q^{2s};q^{-2})_{p-1-k}(q^{2t};q^{-2})_{p-1-k}\\
            &\left(q^{2p-2k-2}(1-q^{2(r+s-p+1)})(1-q^{2(t-p-k-1)})\right.\\
            &\left.\left.-q^{2s}(1-q^{2(l+t-p+1)})(1-q^{2r-2k})\right)\right)\\
        \end{split}
    \end{equation}
    Direct computation shows that
    \begin{equation}\label{eq:4.79}
        \begin{split}
            &q^{2p-2k-2}(1-q^{2(r+s-p+1)})(1-q^{2(t-p-k-1)})\\
            &-q^{2s}(1-q^{2(l+t-p+1)})(1-q^{2r-2k})\\
            =&q^{2p-2k-2}(1-q^{2(s-p-k-1)})(1-q^{2(t-p-k-1)})\\
            &-q^{2t+2s-2(p-1-k)}(1-q^{2(l-k)})(1-q^{2r-2k})
        \end{split}
    \end{equation}
    Substituting (\ref{eq:4.79}) into (\ref{eq:4.78}), we get
    \begin{equation}
        \begin{split}
            &h\left(a^sb^{m-s-r}c^rd^{m-s-l}e^{p}f^{m-r-t}g^lh^{m-l-t}k^t\cdot\text{det}_q^{-m}\right)\\
            % =&\frac{(q^2;q^2)_{r+t-p}(q^2;q^2)_{l+t-p}(q^2;q^2)_{r+s-p}(q^2;q^2)_{l+s-p}}{(q^2;q^2)_m^2}\\
            % &(-1)^{r+l+p}q^{(2m+1)(l+r)+(p-2)m-2l^2-2r^2-rl+4st-3ps-3pt}\\
            % &\left(\sum_{k=0}^{p-1}(-1)^{k}q^{(p-1-k)(p-3k)+2k(s+t)}\right.\\
            % &\cdot{p-1\choose k}_{q^{2}}(q^{2l};q^{-2})_k(q^{2r};q^{-2})_k(q^{2s};q^{-2})_{p-k}(q^{2t};q^{-2})_{p-k}\\
            % &-\sum_{k=0}^{p-1}(-1)^{k}q^{(p-1-k)(p-3k-4)+2(k+1)(s+t)}\\
            % &\left.\cdot{p-1\choose k}_{q^{2}}(q^{2l};q^{-2})_{k+1}(q^{2r};q^{-2})_{k+1}(q^{2s};q^{-2})_{p-1-k}(q^{2t};q^{-2})_{p-1-k}\right)\\
            % =&\frac{(q^2;q^2)_{r+t-p}(q^2;q^2)_{l+t-p}(q^2;q^2)_{r+s-p}(q^2;q^2)_{l+s-p}}{(q^2;q^2)_m^2}\\
            % &(-1)^{r+l+p}q^{(2m+1)(l+r)+(p-2)m-2l^2-2r^2-rl+4st-3ps-3pt}\\
            % &\left(\sum_{k=0}^{p-1}(-1)^{k}q^{(p-k)(p-3k-1)+2k+2k(s+t)}\right.\\
            % &\cdot{p-1\choose k}_{q^{2}}(q^{2l};q^{-2})_k(q^{2r};q^{-2})_k(q^{2s};q^{-2})_{p-k}(q^{2t};q^{-2})_{p-k}\\
            % &+\sum_{k=1}^{p}(-1)^{k}q^{(p-k)(p-3k-1)+2k(s+t)}\\
            % &\left.\cdot{p-1\choose k-1}_{q^{2}}(q^{2l};q^{-2})_{k}(q^{2r};q^{-2})_{k}(q^{2s};q^{-2})_{p-k}(q^{2t};q^{-2})_{p-k}\right)\\
            =&\frac{(q^2;q^2)_{r+t-p}(q^2;q^2)_{l+t-p}(q^2;q^2)_{r+s-p}(q^2;q^2)_{l+s-p}}{(q^2;q^2)_m^2}\\
            &(-1)^{r+l+p}q^{(2m+1)(l+r)+(p-2)m-2l^2-2r^2-rl+4st-3ps-3pt}\\
            &\left(\sum_{k=0}^{p}(-1)^{k}q^{(p-k)(p-3k-1)+2k(s+t)}\left(q^{2k}{p-1\choose k}_{q^{2}}+{p-1\choose k-1}_{q^{2}}\right)\right.\\
            &\left.\cdot(q^{2l};q^{-2})_k(q^{2r};q^{-2})_k(q^{2s};q^{-2})_{p-k}(q^{2t};q^{-2})_{p-k}\right)\\
            =&\frac{(q^2;q^2)_{r+t-p}(q^2;q^2)_{l+t-p}(q^2;q^2)_{r+s-p}(q^2;q^2)_{l+s-p}}{(q^2;q^2)_m^2}\\
            &(-1)^{r+l+p}q^{(2m+1)(l+r)+(p-2)m-2l^2-2r^2-rl+4st-3ps-3pt}\\
            &\left(\sum_{k=0}^{p}(-1)^{k}q^{(p-k)(p-3k-1)+2k(s+t)}\right.\\
            &\left.\cdot{p\choose k}_{q^{2}}(q^{2l};q^{-2})_k(q^{2r};q^{-2})_k(q^{2s};q^{-2})_{p-k}(q^{2t};q^{-2})_{p-k}\right)\\
        \end{split}
    \end{equation}
    The expression is consistent with (\ref{eq:4.77}). Hence, the formula (\ref{eq:4.73}) is true for all possible $n$ values.
\end{proof}

\end{document}
